\documentclass[11pt]{article}
\usepackage[a4paper,margin=25mm]{geometry}
\usepackage{amsmath,amssymb,amsthm}
\usepackage{tikz-cd}
\usepackage{hyperref}

\newtheorem{definition}{Definition}
\newtheorem{proposition}{Proposition}
\newtheorem{remark}{Remark}

\title{Hiroshima Responsibility Functor}
\author{}
\date{}

\begin{document}
\maketitle

\section*{Claim}

Hiroshima placement can be modeled as a functor
\[
  F_H:\mathcal{C}\to\mathcal{R}
\]
from the ordinary corporate decision category to the Hiroshima responsibility
category.  The functor maps the ordinary composite
\[
  D\to P\to O
\]
to the responsibility composite
\[
  D_H\to P_H\to B\to E\to V\to A\to O',
\]
which factors through both Beacon $B$ and Verification $V$.

The theorem-level claim is not that psychological pressure is directly proved.
The theorem-level claim is that paths bypassing $B$ or $V$ are not admissible
under the predicate $\mathrm{Adm}_{B,V}$.

\section*{Why the Subscript $H$ Is Hiroshima}

The subscript $H$ is not a decorative symbol.  It is justified by an external
historical premise: Hiroshima is the place where a nuclear-layer failure of
admissibility is preserved as institutional world memory and is now connected
to international AI governance.

In this formalization, ``nuclear-layer'' means the technological layer at which
humanity can damage the conditions of its own continued existence.  AI is
placed in the same layer not because it is identical to nuclear technology, but
because advanced AI governance also concerns decisions whose failure modes can
scale beyond ordinary local harm.

The Hiroshima-specific premise is the conjunction of three conditions:
\[
\begin{array}{ll}
\text{(i)} & \text{a historical witness of } \neg\mathrm{Adm}_{B,V}
             \text{ at civilizational scale},\\
\text{(ii)}& \text{institutional preservation as world memory},\\
\text{(iii)}& \text{an existing AI governance anchor.}
\end{array}
\]
Condition (i) refers to the atomic bombing of Hiroshima.  Condition (ii) is
represented by institutions such as the Hiroshima Peace Memorial / Genbaku Dome
and the Hiroshima National Peace Memorial Hall for the Atomic Bomb
Victims.\footnote{See UNESCO, ``Hiroshima Peace Memorial (Genbaku Dome)'',
\url{https://whc.unesco.org/en/list/775/}; Hiroshima National Peace Memorial
Hall for the Atomic Bomb Victims,
\url{https://www.hiro-tsuitokinenkan.go.jp/en/about/index.html}.}
Condition (iii) is represented by the Hiroshima AI Process launched under
Japan's G7 presidency in 2023.\footnote{See OECD, ``G7 Hiroshima Process on
Generative Artificial Intelligence (AI)'',
\url{https://www.oecd.org/en/publications/g7-hiroshima-process-on-generative-artificial-intelligence-ai_bf3c0c60-en.html};
JapanGov, ``The Hiroshima AI Process'',
\url{https://www.japan.go.jp/kizuna/2024/02/hiroshima_ai_process.html}.}

This is a structural uniqueness claim about the conjunction of these three
conditions.  It is not a moral exclusion of other cities with histories of war
or technological harm.

\section*{Definitions}

\begin{definition}[Ordinary corporate decision category]
Let $\mathcal{C}$ be the free category generated by
\[
  D \xrightarrow{p} P \xrightarrow{q} O.
\]
Here $D$ is a corporate decision, $P$ is profit/efficiency evaluation, and $O$
is ordinary operational outcome.  The ordinary decision composite is
\[
  u:=q\circ p:D\to O.
\]
\end{definition}

\begin{definition}[Hiroshima responsibility category]
Let $\mathcal{R}$ be the free category generated by
\[
  D_H \xrightarrow{m} P_H \xrightarrow{b} B
  \xrightarrow{e} E \xrightarrow{v} V \xrightarrow{a} A
  \xrightarrow{o} O'.
\]
Here $D_H$ is a Hiroshima-contextualized corporate decision, $P_H$ is a
Hiroshima-contextualized profit/efficiency evaluation, $B$ is the Beacon
admissibility condition, $E$ is evidence, $V$ is verifiable evidence or
third-party verification, $A$ is accountable decision, and $O'$ is accountable
operational outcome.
\end{definition}

\section*{What $D_H$ Does Not Formally Define}

The notation $D_H$ denotes a corporate decision under Hiroshima place-context.
This formalization does not define ``Hiroshima place-context'' internally as a
mathematical object.  Instead, it records the following external premise:

\begin{itemize}
  \item A historical instance structurally interpreted as
  $\neg\mathrm{Adm}_{B,V}$ at civilizational scale exists and is located in
  Hiroshima.
  \item This instance is preserved as institutional world memory.
  \item Hiroshima has already been used as an international AI governance
  reference point through the Hiroshima AI Process.
\end{itemize}

Thus $D_H$ is not arbitrary contextualization.  It is the starting object for
decisions interpreted under a historically documented responsibility context.

\begin{definition}[Hiroshima responsibility functor]
Define the Hiroshima Responsibility Functor
\[
  F_H:\mathcal{C}\to\mathcal{R}
\]
on objects by
\[
  F_H(D)=D_H,\qquad F_H(P)=P_H,\qquad F_H(O)=O'.
\]
Define it on generating morphisms by
\[
  F_H(p)=m:D_H\to P_H,
\]
and
\[
  F_H(q)=o\circ a\circ v\circ e\circ b:P_H\to O'.
\]
Since $\mathcal{C}$ is free on the generators $p,q$, this uniquely extends to
a functor on all morphisms of $\mathcal{C}$.
\end{definition}

\begin{definition}[Admissibility through Beacon and Verification]
For a morphism $f:X\to Y$ in $\mathcal{R}$, define
\[
  f\in \mathrm{Adm}_{B,V}(X,Y)
\]
if and only if there exist morphisms
\[
  x:X\to B,\qquad y:B\to V,\qquad z:V\to Y
\]
such that
\[
  f=z\circ y\circ x.
\]
Equivalently, $f$ factors through both Beacon $B$ and Verification $V$.
\end{definition}

\begin{remark}
In the Lean formalization, $\mathrm{Adm}_{B,V}$ is implemented via
\texttt{pathObjects}, which records the sequence of objects visited by a path.
The equivalence between this implementation and the factorization definition
above is proved as \texttt{AdmBV\_iff\_factorizes} in the Lean file.  The key
auxiliary is a path-splitting lemma: if an object appears in
\texttt{pathObjects}$\,p$, then $p$ splits at that object.
\end{remark}

\section*{Propositions}

\begin{proposition}[Functorial image of ordinary decision]
The Hiroshima responsibility functor sends the ordinary decision composite
$u=q\circ p$ to the responsibility composite
\[
  F_H(u)
  =
  o\circ a\circ v\circ e\circ b\circ m
  :
  D_H\to O'.
\]
\end{proposition}

\begin{proof}
By functoriality,
\[
  F_H(q\circ p)=F_H(q)\circ F_H(p).
\]
Substituting the definitions gives
\[
  F_H(q)\circ F_H(p)
  =
  (o\circ a\circ v\circ e\circ b)\circ m.
\]
By associativity of composition, this is
\[
  o\circ a\circ v\circ e\circ b\circ m:D_H\to O'.
\]
\end{proof}

\begin{proposition}[Responsibility composite is admissible]
The morphism
\[
  o\circ a\circ v\circ e\circ b\circ m:D_H\to O'
\]
belongs to $\mathrm{Adm}_{B,V}(D_H,O')$.
\end{proposition}

\begin{proof}
Let
\[
  x=b\circ m:D_H\to B,\qquad
  y=v\circ e:B\to V,\qquad
  z=o\circ a:V\to O'.
\]
Then
\[
  z\circ y\circ x
  =
  (o\circ a)\circ(v\circ e)\circ(b\circ m)
  =
  o\circ a\circ v\circ e\circ b\circ m.
\]
Hence the responsibility composite factors through both $B$ and $V$.
\end{proof}

\begin{proposition}[Beacon bypass is inadmissible]
For any path $f:X\to Y$ in $\mathcal{R}$, if $f$ does not factor through $B$,
then $f\notin \mathrm{Adm}_{B,V}(X,Y)$.
\end{proposition}

\begin{proof}
By definition, $f\in\mathrm{Adm}_{B,V}(X,Y)$ only if there exist
$x:X\to B$, $y:B\to V$, and $z:V\to Y$ such that
$f=z\circ y\circ x$.  This is a factorization through $B$.  Therefore, if no
such factorization through $B$ exists, $f$ is not admissible.
\end{proof}

\begin{proposition}[Verification bypass is inadmissible]
For any path $f:X\to Y$ in $\mathcal{R}$, if $f$ does not factor through $V$,
then $f\notin \mathrm{Adm}_{B,V}(X,Y)$.
\end{proposition}

\begin{proof}
By definition, $f\in\mathrm{Adm}_{B,V}(X,Y)$ only if there exist
$x:X\to B$, $y:B\to V$, and $z:V\to Y$ such that
$f=z\circ y\circ x$.  This is a factorization through $V$.  Therefore, if no
such factorization through $V$ exists, $f$ is not admissible.
\end{proof}

\section*{Diagram}

\[
\begin{tikzcd}[column sep=large,row sep=large]
\mathcal{C}:
  D \arrow[r,"p"] & P \arrow[r,"q"] & O
\\[1.5em]
\mathcal{R}:
  D_H \arrow[r,"m"] &
  P_H \arrow[r,"b"] &
  B \arrow[r,"e"] &
  E \arrow[r,"v\;(\mathrm{ADIC})"] &
  V \arrow[r,"a"] &
  A \arrow[r,"o"] &
  O'
\end{tikzcd}
\]

The functor $F_H$ maps $D,P,O$ to $D_H,P_H,O'$ and maps the ordinary evaluation
process to the responsibility process.

\section*{Compact Interpretation}

Hiroshima placement is a functor from the ordinary corporate decision category
$\mathcal{C}$ to the Hiroshima responsibility category $\mathcal{R}$:
\[
  F_H:\mathcal{C}\to\mathcal{R}.
\]
It maps corporate decision $D$, profit/efficiency evaluation $P$, and outcome
$O$ to Hiroshima-contextualized decision $D_H$, Hiroshima-contextualized
profit/efficiency evaluation $P_H$, and accountable outcome $O'$.

It sends the ordinary composite
\[
  D\to P\to O
\]
to the responsibility composite
\[
  D_H\to P_H\to B\to E\to V\to A\to O'.
\]
Therefore, a path that bypasses Beacon $B$ or Verification $V$ is not an
admissible responsible morphism under $\mathrm{Adm}_{B,V}$.

\section*{What Is and Is Not Claimed}

\begin{itemize}
  \item Claimed: $F_H$ is a functor from the ordinary decision category to the
  Hiroshima responsibility category.
  \item Claimed: $F_H(q\circ p)=o\circ a\circ v\circ e\circ b\circ m$.
  \item Claimed: this image factors through both Beacon $B$ and Verification
  $V$.
  \item Claimed: paths bypassing $B$ or $V$ are not in
  $\mathrm{Adm}_{B,V}$.
  \item Claimed as an external historical premise: Hiroshima is the unique
  city satisfying the conjunction of (i) nuclear-layer
  $\neg\mathrm{Adm}_{B,V}$ witness, (ii) institutional preservation as world
  memory, and (iii) existing AI governance anchor through the Hiroshima AI
  Process.
  \item Not claimed: Hiroshima uniquely guarantees ethical action.
  \item Not claimed: other cities with war histories are excluded on moral
  grounds.  They are excluded from this specific role only if they lack the
  conjunction required above.
  \item Not claimed: psychological pressure itself is mathematically proved.
  \item Not claimed: category theory alone proves institutional compliance.
\end{itemize}

\end{document}
